\section{Introduction}
\label{section:introduction}
Phase field models \cite{chen2002phase,steinbach2013phase} were introduced to describe the diffusive nature of interfaces between two phases. 
It is known that the Allen-Cahn equation \cite{allen1979microscopic} serves as a fundamental phase field model. 
In this paper, we consider the Allen-Cahn equation of the form
\begin{align}\label{eq:ac}
      \begin{cases}
    \frac{\partial u}{\partial t}-\epsilon^2\Delta u+f(u)=0, &\qquad (x,t)\in \Omega \times (0,+\infty),\\
    u|_{t=0}=u_0,&\qquad x\in \Omega,
\end{cases}
\end{align}
where the domain $\Omega \subset \mathbb{R}^d$, $d=2,3$. $u=u(x,t)$ is the concentration of a phase in a two-phase system and $u_0$ is the initial data. 
The parameter $\epsilon^2>0$ is constant which governs the typical length scale of an interface in the dynamical evolution \cite{wang1993thermodynamically,wheeler1992phase}. 
The nonlinear term $f(u)$ is the derivative of a potential function $F(u)$. 
\par Though the double-well potential is simpler because of its polynomial form, the Flory-Huggins potential is more realistic from the physical viewpoint \cite{flory1942thermodynamics}. 
The logarithmic Flory-Huggins free energy potential is of the form
\begin{equation}
    F(u)=u\log(u)+(1-u)\log(1-u)+\theta(u-u^2),
\end{equation}
where $\theta>2$ is the energy parameter, and respectively,
\begin{equation}
    f(u)=\log(u)-\log(1-u)+\theta(1-2u).
\end{equation}
Defining the continuous energy function,
\begin{equation}
    E(u)=\int_{\Omega} (F(u)+\frac{\epsilon^2}{2}|\nabla u|^2)\text{d}x.
\end{equation}
we then obtain the following natural energy dissipation property \cite{allen1979microscopic}.
\begin{proposition}[Continuous energy stability]\label{prop:sta}
    The energy of a smooth solution $u$ to \eqref{eq:ac} is dissipated with time.
\end{proposition}
\begin{proof}
    By definition, 
    \begin{equation}
        \begin{aligned}
        \frac{\partial E(u)}{\partial t}
        &=\int_{\Omega}(f(u) - \epsilon^2\Delta u)\frac{\partial u}{\partial t}\text{d}x\\
        &=-\int_{\Omega}|\frac{\partial u}{\partial t}|^2\text{d}x\\
        &\leq 0.
    \end{aligned}
    \end{equation}
\end{proof}
Since the equation plays a key role in phase field models and phase transition simulations, 
theoretical analysis and numerical studies concerning it are particularly important. 
A ideal numerical algorithm for the simulation of Allen-Cahn equation with Flory-Huggins potential should possess the following desirable properties: 
(a) bound preserving property: the logarithmic term confines the solution to $0<u<1$,  so the numerical scheme should preserve this bound; 
(b) discrete energy stability: corresponding to the continuous energy stability \cite{provatas2011phase}. 
In what follows, we review some recent work on numerical simulations of Allen-Cahn equation.
\par In \cite{guillen2013linear,yang2017numerical}, 
the so-called invariant energy quadratization (IEQ) method was proposed. 
This method has been widely applied in numerous numerical simulation scenarios, such as molecular beam epitaxy (MBE) \cite{yang2017numerical} and crystal growth \cite{zhao2017numerical,zhao2018general}. 
The basic idea is to introduce some auxiliary variables (SAV) \cite{shen2018scalar} to transform the original free energy functional into a quadratic form. 
Despite its significant success in numerical simulations, the method has two limitations: 
(a) the scheme ensures stability only for the transformed discrete energy, but may not for the discrete energy of the original formulation 
\cite{yang2017numerical,wang2020stabilized}; 
(b) the time step size required to satisfy the maximum principle \cite{chen1992generation,nishiura1987stability} is restricted by a uniform constant, 
which introduces a trade-off between stability and computational efficiency.
\par Reference \cite{du2019maximum} introduced the first order exponential time differencing (ETD1) method for the nonlocal Allen-Cahn equation \cite{kroemer2025quantitative}. 
The scheme is shown to preserve the maximum principle and exhibit asymptotic compatibility. 
Unfortunately, this article only investigates the Allen-Cahn equation with the double-well potential. 
Upon reflection, we find that for the Allen-Cahn equation with Flory-Huggins potential, the algorithm faces limitations, 
because the validity of the discrete maximum principle requires the potential term to be bounded, 
a property that the Flory–Huggins potential does not satisfy. 
Furthermore, the maximum principle is only theoretically guaranteed, which means, 
in numerical practice, we need the solution strictly satisfy the separation property, i.e. $\|u\|_{\infty}\leq 1-\delta$ for sufficiently small $\delta$.
\par \cite{li2022stability} proposed an algorithm based on operator splitting methods \cite{macnamara2017operator}. 
It applied the Strang splitting discretization \cite{strang1968construction} for second-order in time splitting methods. 
The merits of the proposed algorithm include its simplicity, its preservation of the maximum principle and discrete energy dissipation, as well as its proven convergence. 
Nonetheless, the method exhibits a shortcoming: the maximum principle is contingent upon a time step restriction. 
Consequently, this imposes challenges for long-time simulations and for computations employing large time steps.
\par A recent work \cite{wang2020stabilized} introduced the stabilized energy factorization (SEF) approach. 
Its core idea is to factorize the energy functional into a product of several factors  \cite{kou2020novel}, 
with each factor handled individually based on its unique property. 
This approach offers multiple benefits: linear semi-implicit scheme, preservation of the discrete maximum principle and discrete energy stability, as well as convergence property. 
However, in the paper, the stability parameter condition required for the maximum principle is stringent and difficult to choose in practical computations \cite{wang2020stabilized}. 
Moreover, the algorithm struggles to fully overcome the numerical instability caused by the singularity at $u=0,1$.
\par The convex splitting scheme was initially proposed in the study of unconditionally energy stable schemes for the Cahn-Hilliard equation by Eyre \cite{eyre1998unconditionally}, 
and was subsequently successfully extended and applied to various phase field models \cite{chen2012linear,shen2012second,baskaran2013convergence}. 
The fundamental principle of the convex splitting approach is to express the energy functional as the sum of a convex and a concave contribution. 
Correspondingly, the convex part is handled implicitly and the concave part explicitly, 
leading to a semi-implicit scheme known as a convex splitting scheme, 
which possesses several desirable properties, such as unique solvability and discrete energy stability \cite{eyre1998unconditionally,du2020phase}. 
Applying this method to the Allen-Cahn equation with Flory-Huggins potential results in a nonlinear discrete system at each time step, 
which necessitates iterative solvers. 
Solving such a system is complex and difficult, primarily due to the strong nonlinearity of the logarithmic potential term and its singularity at $u=0,1$.
\par THERE NEEDS SOME MATERIALS ABOUT ITERATIVE SOLVER.
\par In this paper, we present a novel iterative solver for the nonlinear system at each time step arising from the application of the convex splitting scheme to the Allen-Cahn equation with Flory–Huggins potential. 
We reformulate this nonlinear system equivalently as a convex optimization problem for finding a minimizer. 
To solve this optimization problem, we employ the alternating direction method of multipliers (ADMM) framework \cite{neal2011distributed}. 
More specifically, the objective functional is decomposed into a linear part and a logarithmic part. 
The subproblem corresponding to the linear part is solved efficiently with the fast Fourier transform (FFT) techniques \cite{cooley1965algorithm}, due to its coefficient matrix. 
For the logarithmic part, the corresponding subproblem is reduced to solving multiple one-dimensional equations with a monotonic derivative, 
which enables us to apply the Newton's method to each equation \cite{dennis1996numerical,ortega2000iterative} in parallel.
\par The remainder of this paper is organized as follows. 
In \Cref{section:discrete}, we present the problem formulation and introduce the discrete function spaces and operators within the framework of the cell-centered finite difference (CCFD) method in spatial discretization. 
\Cref{section:admm} details the convex splitting scheme for the equation and our novel iterative solver; 
for comparative purposes, we also describe the conventional solver based on high-dimensional Newton's method. 
\Cref{section:analysis} provides a theoretical analysis of the scheme, addressing the well-posedness and discrete energy stability of the convex splitting scheme, 
as well as the unconditional convergence and bound-preserving property of the solver. 
Numerical experiments that illustrate the effectiveness and robustness of the proposed solver are reported in \Cref{sec:numerical}. 
Finally, \Cref{section:conclusion} offers some concluding remarks.

